% !TEX program = pdflatex
\documentclass[11pt,a4paper]{article}

% =========================================================================
% PLAN DE INFORMACIÓN Y COMUNICACIONES
% Proyecto: CasaScore (AlquilaPuno)
% Gestión de Proyectos TICS (EST315) -- Semestre 2026-II -- Team Sobrevivientes
% =========================================================================

% --- PAQUETES DE CONFIGURACIÓN ---
\usepackage[utf8]{inputenc}
\usepackage[T1]{fontenc}
\usepackage[spanish,es-tabla]{babel}
% Compatibilidad con babel-spanish reciente (flechas de TikZ y signo de porcentaje)
\AtBeginDocument{\shorthandoff{<>}\spanishplainpercent}
\usepackage{geometry}
\geometry{
    top=3.6cm,
    bottom=2.3cm,
    left=2.2cm,
    right=2.2cm,
    % PARTE MODIFICABLE: headheight debe ser mayor que la altura del logo más alto (2.56cm = 73pt).
    headheight=80pt,
    headsep=15pt,
    footskip=20pt
}

\usepackage{graphicx}
\usepackage{fancyhdr}
\usepackage[table]{xcolor}
\usepackage{titlesec}
\usepackage{booktabs}
\usepackage{tabularx}
\usepackage{array}
\usepackage{amsmath}
\usepackage{enumitem}
\usepackage{microtype}
\usepackage{amssymb}
% --- RECURSOS VISUALES: GRÁFICOS Y CÓDIGO DE COLORES ---
\usepackage{tikz}
\usetikzlibrary{positioning,calc}
\definecolor{GrisFila}{RGB}{246, 248, 251}
\definecolor{CuadCerca}{RGB}{150, 30, 30}
\definecolor{CuadSatisf}{RGB}{180, 130, 20}
\definecolor{CuadInfo}{RGB}{10, 45, 90}
\definecolor{CuadMonit}{RGB}{110, 115, 125}
\newcommand{\etiqueta}[2]{{\setlength{\fboxsep}{2pt}\colorbox{#1}{\textcolor{white}{\scriptsize\bfseries #2}}}}
\newcommand{\cerca}{\etiqueta{CuadCerca}{Gestionar de cerca}}
\newcommand{\satisf}{\etiqueta{CuadSatisf}{Mantener satisfecho}}
\newcommand{\informado}{\etiqueta{CuadInfo}{Mantener informado}}
\newcommand{\monit}{\etiqueta{CuadMonit}{Monitorear}}
\usepackage{needspace}
\usepackage{longtable}
\definecolor{Rama1}{RGB}{10, 45, 90}
\definecolor{Rama2}{RGB}{40, 95, 150}
\definecolor{Rama3}{RGB}{180, 130, 20}
\definecolor{Rama4}{RGB}{30, 110, 60}
\definecolor{Rama5}{RGB}{110, 60, 140}
\definecolor{RaciR}{RGB}{150, 30, 30}
\definecolor{RaciA}{RGB}{10, 45, 90}
\definecolor{RaciC}{RGB}{180, 130, 20}
\definecolor{RaciI}{RGB}{140, 145, 155}
\definecolor{NivAlto}{RGB}{30, 110, 60}
\definecolor{NivMedio}{RGB}{180, 130, 20}
\definecolor{NivBajo}{RGB}{150, 30, 30}
\usetikzlibrary{patterns}
\definecolor{SevBaja}{RGB}{30, 110, 60}
\definecolor{SevMedia}{RGB}{180, 130, 20}
\definecolor{SevAlta}{RGB}{150, 30, 30}
\usetikzlibrary{arrows.meta,shapes.geometric,fit,backgrounds}
\usepackage{listings}
\lstset{basicstyle=\ttfamily\scriptsize, frame=single, rulecolor=\color{AzulUNA!40}, backgroundcolor=\color{GrisFila}, columns=fullflexible, keepspaces=true, breaklines=true, xleftmargin=4pt, xrightmargin=4pt, literate={á}{{\'a}}1 {é}{{\'e}}1 {í}{{\'i}}1 {ó}{{\'o}}1 {ú}{{\'u}}1 {ñ}{{\~n}}1 {Á}{{\'A}}1}

\usepackage{hyperref}

% Logos: funciona a 2 o 3 niveles de profundidad dentro de docs/latex/
\graphicspath{{../../assets/img/}{../../../assets/img/}}

% --- PALETA INSTITUCIONAL FIEI / UNA PUNO ---
\definecolor{AzulUNA}{RGB}{10, 45, 90}
\definecolor{DoradoFIEI}{RGB}{180, 130, 20}
\definecolor{GrisOscuro}{RGB}{50, 50, 50}
\definecolor{AzulClaro}{RGB}{228, 236, 246}

% --- CONFIGURACIÓN DE LISTAS ---
\setlist{topsep=2pt, itemsep=2pt, parsep=1pt}

% --- COLUMNAS DE TABLA ---
\newcolumntype{L}{>{\raggedright\arraybackslash}X}
\newcolumntype{P}[1]{>{\raggedright\arraybackslash}p{#1}}
\newcolumntype{C}[1]{>{\centering\arraybackslash}p{#1}}
\newcolumntype{Q}[1]{>{\raggedleft\arraybackslash}p{#1}}
\newcolumntype{R}{>{\raggedleft\arraybackslash}X}
\renewcommand{\arraystretch}{1.25}

% --- FORMATO DE ENLACES ---
\hypersetup{
    colorlinks=true,
    linkcolor=AzulUNA,
    urlcolor=AzulUNA,
    citecolor=AzulUNA,
    pdftitle={Plan de Información y Comunicaciones -- CasaScore (AlquilaPuno)},
    pdfauthor={Team Sobrevivientes}
}

% --- FORMATO DE TÍTULOS ---
\titleformat{\section}{\color{AzulUNA}\normalfont\large\bfseries}{\thesection.}{0.7em}{}[\color{AzulUNA}\titlerule]
\titleformat{\subsection}{\color{AzulUNA}\normalfont\normalsize\bfseries}{\thesubsection.}{0.5em}{}
\titlespacing*{\section}{0pt}{9pt}{4pt}
\titlespacing*{\subsection}{0pt}{6pt}{2pt}

% --- ENCABEZADO Y PIE DE PÁGINA (fancyhdr) ---
\pagestyle{fancy}
\fancyhf{}

% ===== TAMAÑO DE LOS LOGOS (PARTE MODIFICABLE) =====
% logo_fiei.png  : height=1.45cm (lado izquierdo)
% logo_equipo.png: height=2.56cm (lado derecho; 25 % más grande porque la imagen tiene más borde blanco)
% Para cambiar el tamaño, modifique solo el valor height=... de cada \includegraphics.
% Si el logo más alto supera 2.6cm, aumente también headheight y top en la configuración de página.
\fancyhead[L]{%
    \includegraphics[height=1.45cm,keepaspectratio]{logo_fiei.png}%
}
\fancyhead[C]{%
    \scriptsize \textbf{\textcolor{AzulUNA}{UNIVERSIDAD NACIONAL DEL ALTIPLANO DE PUNO}}\\[2pt]
    \tiny \textcolor{GrisOscuro}{FACULTAD DE INGENIERÍA ESTADÍSTICA E INFORMÁTICA}\\[2pt]
    \scriptsize \textbf{\textcolor{AzulUNA}{Gestión de Proyectos TICS (EST315) \textbar\ Team Sobrevivientes}}\\[2pt]
    \footnotesize \textbf{\textcolor{DoradoFIEI}{PLAN DE INFORMACIÓN Y COMUNICACIONES}}%
}
\fancyhead[R]{%
    \includegraphics[height=2.56cm,keepaspectratio]{logo_equipo.png}%
}

\renewcommand{\headrulewidth}{0pt}
\renewcommand{\headrule}{}

\fancyfoot[L]{\scriptsize \textcolor{GrisOscuro}{CasaScore (AlquilaPuno)}}
\fancyfoot[C]{\small \thepage}
\fancyfoot[R]{\scriptsize \textcolor{GrisOscuro}{Puno, Perú -- Semestre 2026-II}}
\renewcommand{\footrulewidth}{0.4pt}

\begin{document}

\begin{center}
    {\Large \textbf{\textcolor{AzulUNA}{PLAN DE INFORMACIÓN Y COMUNICACIONES}}}\\[3pt]
    {\normalsize \textbf{CasaScore (AlquilaPuno)}}\\[4pt]
    \footnotesize \textbf{Fecha:} martes 22/09/2026 (semana 5) \quad \textbar\quad \textbf{Versión:} 1.0 \quad \textbar\quad \textbf{Fase:} Planificación
\end{center}

\vspace{0.1cm}

\section{Datos Generales}
\noindent\begin{tabularx}{\textwidth}{@{}P{4.2cm} L@{}}
\toprule
\textbf{Proyecto:} & CasaScore (aplicación AlquilaPuno) \\
\textbf{Responsable:} & Viviana Mary Quisbert Quispe -- Gestora de Comunicación, Riesgos y Documentación \\
\textbf{Proceso PM4DEV:} & Gestión de la información y las comunicaciones \\
\textbf{Insumos:} & Registro de interesados, plan del equipo, cronograma y plan de riesgos \\
\bottomrule
\end{tabularx}

\section{Objetivo}
Asegurar que cada interesado reciba la información que necesita, en el momento oportuno y por el canal adecuado, y que toda decisión del proyecto quede registrada.

\section{Flujo de Información}
\begin{center}
\begin{tikzpicture}[font=\scriptsize, caja/.style={draw=#1, fill=#1!8, line width=0.8pt, rounded corners=3pt, text width=3.3cm, align=center, inner sep=4pt, minimum height=1.15cm}, >=Stealth]
\node[caja=AzulUNA, text width=3.4cm, minimum height=1.7cm] (eq) at (0,0) {\textbf{\textcolor{AzulUNA}{TEAM SOBREVIVIENTES}}\\[2pt] WhatsApp, GitHub y Meet\\[1pt] {\tiny actas, informes y tablero}};
\node[caja=Rama4] (ar) at (-5.3,1.25) {\textbf{\textcolor{Rama4}{Arrendatarios}}\\ {\tiny grupos de WhatsApp, afiches y encuestas}};
\node[caja=Rama3] (pr) at (-5.3,-1.25) {\textbf{\textcolor{Rama3}{Propietarios}}\\ {\tiny visitas, guía y WhatsApp}};
\node[caja=Rama2] (as) at (5.3,1.25) {\textbf{\textcolor{Rama2}{Docente asesor}}\\ {\tiny entregables en la sesión de clase}};
\node[caja=GrisOscuro] (ot) at (5.3,-1.25) {\textbf{\textcolor{GrisOscuro}{Concursos, anunciantes y proveedores}}\\ {\tiny según la matriz}};
\foreach \n/\c in {ar/Rama4, pr/Rama3, as/Rama2, ot/GrisOscuro} \draw[<->, thick, \c] (eq) -- (\n);
\end{tikzpicture}
\end{center}

\section{Canales de Comunicación}
{\small
\rowcolors{2}{white}{GrisFila}
\noindent\begin{tabularx}{\textwidth}{@{}P{3.2cm} C{1.6cm} L P{2.6cm}@{}}
\toprule
\rowcolor{AzulClaro}
\textbf{Canal} & \textbf{Uso} & \textbf{Para qué se usa} & \textbf{Regla} \\
\midrule
Grupo de WhatsApp del equipo & \etiqueta{AzulUNA}{Interno} & Avisos, coordinación diaria y alertas de riesgos. & Respuesta en 12 h o menos \\
GitHub (Issues y Projects) & \etiqueta{AzulUNA}{Interno} & Tareas, avance del código, actas y documentos versionados. & Toda tarea es un Issue \\
Google Meet & \etiqueta{AzulUNA}{Interno} & Reuniones de control y de revisión del sprint. & Cámara y agenda previa \\
Grupos de WhatsApp de clientes & \etiqueta{Rama4}{Externo} & Convocatorias, avances, enlaces de prueba y encuestas. & Mensajes formales y breves \\
Afiches y visitas & \etiqueta{Rama4}{Externo} & Difusión en la ciudad universitaria y trato directo con propietarios. & Siempre dos integrantes \\
Sesión de clase & \etiqueta{Rama2}{Externo} & Presentación de entregables y consultas al docente asesor. & Según la ruta del curso \\
\bottomrule
\end{tabularx}
}

\Needspace{12\baselineskip}
\section{Necesidades de Información}
{\small
\rowcolors{2}{white}{GrisFila}
\noindent\begin{tabularx}{\textwidth}{@{}P{3.6cm} L P{3.4cm}@{}}
\toprule
\rowcolor{AzulClaro}
\textbf{Interesado} & \textbf{Qué necesita saber} & \textbf{Para qué} \\
\midrule
Team Sobrevivientes & Avance, gastos, riesgos, cambios y tareas de cada integrante. & Decidir como sponsor \\
Arrendatarios & Cómo se calcula el puntaje, cuándo es el piloto y cómo usar la app. & Confiar en el ranking \\
Propietarios & Qué datos y fotos se piden, cómo publicar y qué beneficio obtienen. & Publicar sus cuartos \\
Docente asesor & Entregables del curso y su aplicación de PM4DEV. & Orientar al equipo \\
Concursos y anunciantes & Propuesta de valor, resultados del piloto y alcance. & Evaluar o invertir \\
Proveedores & Uso del servicio dentro de sus términos. & Mantener la cuenta activa \\
\bottomrule
\end{tabularx}
}

\section{Matriz de Comunicaciones}
{\footnotesize
\setlength{\tabcolsep}{3.5pt}
\begin{longtable}{@{}>{\raggedright\arraybackslash}p{4.0cm} >{\raggedright\arraybackslash}p{2.9cm} >{\raggedright\arraybackslash}p{2.4cm} >{\raggedright\arraybackslash}p{2.9cm} >{\raggedright\arraybackslash}p{2.7cm}@{}}
\toprule
\rowcolor{AzulClaro}
\textbf{Qué información} & \textbf{Para quién} & \textbf{Cuándo} & \textbf{Canal} & \textbf{Quién la produce} \\
\midrule
\endfirsthead
\toprule
\rowcolor{AzulClaro}
\textbf{Qué información} & \textbf{Para quién} & \textbf{Cuándo} & \textbf{Canal} & \textbf{Quién la produce} \\
\midrule
\endhead
\bottomrule
\endfoot
Acta de la reunión de control & Equipo & Martes, dentro de las 24 h & GitHub (carpeta de actas) & Viviana M. Quisbert Quispe \\
\rowcolor{GrisFila}
Informe de avance: valor ganado, SPI, CPI y riesgos & Equipo & Martes & Google Meet y GitHub & Ruth K. Apaza Solis \\
Revisión del sprint y demostración del software & Equipo & Viernes & Google Meet & Eliseo Tarqui Ajahuana \\
\rowcolor{GrisFila}
Tablero de tareas y responsables & Equipo & Siempre actualizado & GitHub Projects & Clyde N. Paricahua Pari \\
Avisos urgentes y riesgos activados & Equipo & Cuando ocurren & Grupo de WhatsApp & Quien lo detecta \\
\rowcolor{GrisFila}
Entregables del curso y sus observaciones & Docente asesor & Según la ruta de entregables & Sesión de clase & Clyde N. Paricahua Pari \\
Convocatoria a entrevistas y al piloto & Arrendatarios estudiantes & Semanas 6, 12 y 13 & Grupos de WhatsApp y afiches & Viviana M. Quisbert Quispe \\
\rowcolor{GrisFila}
Guía para publicar y acompañamiento & Propietarios & Semanas 6 y 13 & Visitas y WhatsApp & Clyde N. Paricahua Pari \\
Encuesta de satisfacción & Arrendatarios y propietarios del piloto & Semana 13 & Google Forms enviado por WhatsApp & Ruth K. Apaza Solis \\
\rowcolor{GrisFila}
Avances e invitación a probar la app & Arrendatarios trabajadores y familias & Semanas 13 y 16 & Grupos de WhatsApp & Viviana M. Quisbert Quispe \\
Postulación con el MVP & Organizadores de concursos y ferias & Semanas 16 y 17 & Medio que fijen las bases & Ruth K. Apaza Solis \\
\rowcolor{GrisFila}
Lista de negocios para el año 1 & Anunciantes locales & Semana 17 & Visitas & Ruth K. Apaza Solis \\
Revisión de términos y límites & Proveedores de servicios & Mensual & Portal de cada proveedor & Eliseo Tarqui Ajahuana \\
\end{longtable}
}

\Needspace{12\baselineskip}
\section{Reuniones}
{\small
\rowcolors{2}{white}{GrisFila}
\noindent\begin{tabularx}{\textwidth}{@{}P{2.7cm} P{2.6cm} L P{2.6cm}@{}}
\toprule
\rowcolor{AzulClaro}
\textbf{Reunión} & \textbf{Día y hora} & \textbf{Agenda} & \textbf{Resultado} \\
\midrule
Control & Martes, 20:00 a 21:00 & Avance contra la línea base, valor ganado, riesgos y cambios. & Acta e informe de avance \\
Revisión del sprint & Viernes, 20:00 a 21:00 & Demostración, retrospectiva y planificación del siguiente sprint. & Acta del sprint \\
Extraordinaria & Cuando se activa un riesgo alto & Análisis del problema y respuesta. & Acta con acuerdos \\
\bottomrule
\end{tabularx}
}

\vspace{0.1cm}
\begin{itemize}[leftmargin=*]
    \item La agenda se publica en el grupo de WhatsApp al menos 12 horas antes.
    \item El acta registra asistentes, acuerdos, responsables y fechas, y se sube a GitHub dentro de las 24 horas.
\end{itemize}

\section{Gestión de Documentos}
\begin{itemize}[leftmargin=*]
    \item \textbf{Repositorio único:} todos los documentos y el código se guardan en GitHub; la versión vigente es la de la rama principal.
    \item \textbf{Nombres:} número del entregable, nombre corto y proyecto, por ejemplo, \texttt{10\_plan\_riesgos\_} \texttt{casascore.tex}.
    \item \textbf{Versiones:} 1.0 al aprobarse; 1.1, 1.2 por correcciones; 2.0 por cambios aprobados con una solicitud de cambio.
    \item \textbf{Datos personales:} las encuestas y los datos de los clientes no se publican en los grupos ni en el repositorio público.
\end{itemize}

\Needspace{12\baselineskip}
\section{Escalamiento de Problemas}
\begin{center}
\begin{tikzpicture}[font=\scriptsize, nv/.style={draw=#1, fill=#1!10, rounded corners=3pt, text width=3.0cm, align=center, minimum height=1.15cm, line width=0.8pt}, >=Stealth]
\node[nv=Rama4] (a) at (0,0) {\textbf{Nivel 1}\\ Responsable de la tarea\\ {\tiny resuelve en 24 h}};
\node[nv=SevMedia] (b) at (5.0,0) {\textbf{Nivel 2}\\ Project Manager\\ {\tiny decide en la siguiente reunión}};
\node[nv=SevAlta] (c) at (10.0,0) {\textbf{Nivel 3}\\ Equipo como sponsor\\ {\tiny decide por unanimidad}};
\draw[->, thick, GrisOscuro] (a) -- node[above, font=\tiny, text width=1.6cm, align=center]{si no se resuelve} (b);
\draw[->, thick, GrisOscuro] (b) -- node[above, font=\tiny, text width=1.6cm, align=center]{si afecta una línea base} (c);
\end{tikzpicture}
\end{center}
Si el problema requiere orientación metodológica, el Project Manager lo consulta con el docente asesor.

\Needspace{24\baselineskip}
\section{Firmas}
\noindent En señal de conformidad con el presente documento, los integrantes del equipo firman a continuación:

\vspace{1.5cm}
\begin{center}
\begin{tabular}{c c c}
    \parbox{6.5cm}{\centering \rule{6.0cm}{0.6pt}\\[5pt] \textbf{Clyde Neil Paricahua Pari}\\[2pt] {\footnotesize Project Manager}}
    & \hspace{0.8cm} &
    \parbox{6.5cm}{\centering \rule{6.0cm}{0.6pt}\\[5pt] \textbf{Ruth Karina Apaza Solis}\\[2pt] {\footnotesize Analista de Planificación}}
    \\[2.0cm]
    \parbox{6.5cm}{\centering \rule{6.0cm}{0.6pt}\\[5pt] \textbf{Eliseo Tarqui Ajahuana}\\[2pt] {\footnotesize Especialista Técnico / Arquitectura}}
    & \hspace{0.8cm} &
    \parbox{6.5cm}{\centering \rule{6.0cm}{0.6pt}\\[5pt] \textbf{Viviana Mary Quisbert Quispe}\\[2pt] {\footnotesize Gestora de Comunicación, Riesgos y Documentación}}
\end{tabular}
\end{center}
\end{document}
